\begin{center}
\textsc{Gauss} an \textsc{Bessel}.\quad Göttingen, 20.~November 1826.
\end{center}

\par\noindent\hspace*{\parindent}.\ .\ .\ .\ .\ .\quad Die Aufbereitung des Materials für das geplante Werk über meine Messungen kostet mich viel Zeit.
Die Berechnung meiner Hauptdreiecke, die 33 Punkte umfassen, ist zwar längst abgeschlossen, aber die Berechnung der vielen durch Richtungsschnitte erfassten Nebenpunkte, die für die Geografie eines bedeutenden Teils Norddeutschlands wichtig sind, macht viel Arbeit, da diese bisher entweder noch gar nicht oder nur vorläufig berechnet worden waren.
Mein Verzeichnis enthält jetzt etwa 250 Punkte, alle aus dem nördlicheren Teil meiner Messungen.
Da meine Dreiecke über die kurhessischen mit den bayerischen und württembergischen zusammenhängen und \textsc{Bohnenberger} mir letztere freundlicherweise mitgeteilt hat (ebenso wie \textsc{Eckhardt} die darmstädtischen), ist es mir sehr unangenehm, dass alle meine Bemühungen, die bayerischen zu erhalten,
\par\smallskip\noindent{\footnotesize IX.}\hfill{\footnotesize46}\par
\clearpage
\noindent362\hfill{\small\textbf{BRIEFWECHSEL MIT BESSEL.}}\par\medskip
\noindent bisher vergeblich waren. \textsc{Schumacher} erzählte mir, \textsc{Soldner} habe ihm gesagt, die Kommission in München habe meine Bitte um Mitteilung deshalb nicht erfüllt, weil man annehme, dass ich zu diesen Dreiecken Berechnungen anstellen wolle!
Ebenso ist es mir mit den österreichischen ergangen.
Auch General \textsc{von Müffling} sollte doch seine Dreiecke östlich von Seeberg bekannt machen, die Berlin anschließen und sich bis über Schlesien erstrecken.

\par\noindent\hspace*{\parindent}Noch größere Schwierigkeiten bereitet mir der weit umfangreichere theoretische Teil, der so vielfach in andere Teile der Mathematik hineinreicht.
Ich sehe hier keinen anderen Ausweg, als mehrere große Hauptabschnitte aus dem Werk herauszulösen, damit sie eigenständig und in der nötigen Ausführlichkeit entwickelt werden können.
Gewissermaßen habe ich damit schon in meiner Schrift über die Abbildung der Flächen unter Erhaltung der Ähnlichkeit der kleinsten Teile begonnen. Eine zweite Abhandlung, die ich vor ein paar Monaten der königlichen Sozietät übergeben habe und die hoffentlich bald gedruckt wird, enthält die Grundsätze und Methoden zur Ausgleichung der Messungen (nebenbei lässt sich daraus indirekt auch ersehen, wie weit die \textsc{Krayenhoff}schen Messungen von der Genauigkeit entfernt sind, die man ihnen zu Unrecht zugeschrieben hat).
Vielleicht werde ich zunächst noch eine dritte Abhandlung ausarbeiten, die verschiedene neue Lehrsätze über gekrümmte Flächen, kürzeste Linien, die Darstellung gekrümmter Flächen in der Ebene usw.\ entwickeln wird.
Hätte ich all diese Gegenstände in mein geplantes Werk aufnehmen sollen, so hätte ich entweder manches oberflächlich behandeln oder dem Werk ein sehr uneinheitliches Aussehen geben müssen.\ .\ .\ .\ .\ .\ .
